\documentclass[10pt]{article}
\usepackage[paperwidth=16cm,paperheight=9cm,margin=0.65cm]{geometry}
\usepackage[spanish,es-noshorthands]{babel}
\usepackage[utf8]{inputenc}
\usepackage[T1]{fontenc}
\usepackage{lmodern}
\usepackage{amsmath}
\usepackage{booktabs,tabularx,array}
\usepackage{xcolor}
\usepackage{tikz}
\usetikzlibrary{arrows.meta,positioning}
\pagestyle{empty}

\definecolor{UniPrimary}{HTML}{E8552D}
\definecolor{UniPrimaryDeep}{HTML}{C44624}
\definecolor{UniLight}{HTML}{F6F7F9}
\newcolumntype{Y}{>{\raggedright\arraybackslash}X}
\setlength{\parindent}{0pt}
\setlength{\parskip}{0.18cm}
\renewcommand{\familydefault}{\sfdefault}
\newcounter{slide}
\newcommand{\slide}[2]{%
  \ifnum\value{slide}>0\clearpage\fi
  \stepcounter{slide}
  \begin{tikzpicture}[remember picture,overlay]
    \fill[UniPrimaryDeep] (current page.north west) rectangle ([yshift=-1.0cm]current page.north east);
    \node[anchor=west,text=white,font=\bfseries\large] at ([xshift=0.65cm,yshift=-0.52cm]current page.north west) {#1};
    \node[anchor=east,text=white,font=\small] at ([xshift=-0.65cm,yshift=-0.52cm]current page.north east) {\theslide};
  \end{tikzpicture}
  \vspace*{0.72cm}
  #2
}

\begin{document}
\slide{Programacion y control de un robot Pioneer}{
{\Large Actividad 2 - CoppeliaSim/Lua}\\
Universidad Internacional de Valencia\\[0.2cm]
Autor: Jorge Luis Mayorga Taborda\\
Profesor: Jose I. Iniguez\\
Entrega: 26 de mayo de 2026
}

\slide{Objetivo}{
\begin{itemize}
  \item Programar un Pioneer P3DX en CoppeliaSim.
  \item Parte 1: desplazamiento hacia un objetivo mediante campos potenciales.
  \item Parte 2: evitacion reactiva de obstaculos.
  \item Integracion del robot en una celda robotizada.
\end{itemize}
}

\slide{Trazabilidad 5/5}{
\begin{tabularx}{\textwidth}{Y Y}
\toprule
\textbf{Guia y temas} & \textbf{Implementacion final}\\
\midrule
Pioneer en CoppeliaSim/Lua & Script Lua embebido en PioneerP3DX.\\
Campos potenciales & Atraccion a mannequin con saturacion de velocidad.\\
Anticolision & 16 ultrasonidos + 6 sensores virtuales.\\
Celda robotizada & Docking/carga, HMI, transportadores, palets, vallas y eventos.\\
Bateria y robustez & Zonas de velocidad, sensor degradado y telemetria industrial.\\
A* cooperativo & 3 AMR auxiliares con reservas temporales, bateria y recarga.\\
PPTs de clase & Modelo diferencial, sensores, FSM, navegacion y validacion.\\
\bottomrule
\end{tabularx}
}

\slide{Entorno de simulacion}{
\begin{itemize}
  \item CoppeliaSim con scripts Lua embebidos.
  \item Escenas base: Actividad2\_1\_Pioneer.ttt y Actividad2\_2\_Pioneer.ttt.
  \item Escena final: Actividad2\_Pioneer\_Profesional\_10\_10.ttt.
  \item Objetivos posibles: mannequin, Bill o dummy de estacion.
  \item 16 ultrasonidos + 6 sensores virtuales de seguridad.
  \item Escenarios: pallet movil, parada, objetivo dinamico, pasillo estrecho, sensor degradado y bateria baja.
\end{itemize}
}

\slide{Modelo Pioneer P3DX}{
\[
v_L = v - \frac{L}{2}\omega,\qquad
v_R = v + \frac{L}{2}\omega
\]
\[
\omega_L = \frac{v_L}{r},\qquad
\omega_R = \frac{v_R}{r}
\]
Robot diferencial con control por velocidad de ruedas y sensores de proximidad.
}

\slide{Parte 1: campos potenciales}{
\[
\vec{F}_{att}=k_{att}
\begin{bmatrix}
x_g-x_r\\
y_g-y_r
\end{bmatrix}
\]
\begin{itemize}
  \item Calcular vector robot-objetivo.
  \item Obtener rumbo deseado con atan2.
  \item Regular avance y giro con saturacion.
  \item Detener al entrar en tolerancia.
\end{itemize}
}

\slide{Parte 2: anticolisiones}{
\begin{itemize}
  \item Cada sensor cercano genera fuerza repulsiva.
  \item La direccion final combina atraccion y repulsion.
  \item El robot reduce velocidad cuando hay obstaculos cerca.
  \item Se prueban obstaculos frontales, laterales, objetivo movil, ruido y fallo intermitente de sensor.
\end{itemize}
}

\slide{Flujo de control}{
\begin{tikzpicture}[node distance=8mm,box/.style={draw=UniPrimaryDeep,rounded corners,fill=UniLight,align=center,text width=2.7cm,minimum height=10mm},arr/.style={-{Latex},UniPrimaryDeep,thick}]
\node[box] (read) {Leer objetivo, ruta y sensores};
\node[box,right=of read] (force) {FSM + campos potenciales};
\node[box,right=of force] (cmd) {Saturar v, omega};
\node[box,below=of force] (wheels) {Enviar velocidades de rueda};
\draw[arr] (read)--(force);
\draw[arr] (force)--(cmd);
\draw[arr] (cmd)--(wheels);
\draw[arr] (wheels.west)--++(-1.7,0)|-(read.south);
\end{tikzpicture}
}

\slide{Integracion en celda}{
\begin{itemize}
  \item Zona segura de espera.
  \item Mision hacia Bill, mannequin o punto de estacion.
  \item Evitacion de obstaculos de la celda.
  \item Senales internas: missionReady, pioneerArrived, pioneerState y pioneerScenario.
  \item Celda con transportadores, palets, vallas, eventos, estacion objetivo y estacion de carga.
\end{itemize}
}

\slide{Extensiones robustas}{
\begin{itemize}
  \item Bateria simulada con modos NORMAL, LOW y CHARGING.
  \item Zonas de velocidad: docking, carril de aproximacion y zona por defecto.
  \item Filtro de ruido y fallo intermitente en sensores fronto-laterales.
  \item HMI visual con bateria, estado, zona activa, carga y fallo sensorial.
  \item Flota auxiliar de 3 AMR con A* cooperativo y retorno a estaciones de carga.
\end{itemize}
}

\slide{Flota A* cooperativa}{
\begin{itemize}
  \item 3 robots AMR auxiliares sobre una grilla de celda.
  \item A* con reservas celda-tiempo y arista-tiempo para evitar ocupaciones simultaneas y cruces opuestos.
  \item Bateria con descarga, retorno autonomo a estacion de carga y recarga hasta READY.
  \item Resultado validado: 16 planes A*, 8 conflictos evitados, 4 eventos de carga y 3/3 robots recargados.
\end{itemize}
}

\slide{Validacion}{
\begin{tabularx}{\textwidth}{Y Y}
\toprule
\textbf{Prueba} & \textbf{Criterio}\\
\midrule
Atraccion & Distancia al objetivo decrece.\\
Llegada & Parada dentro de tolerancia.\\
Obstaculo frontal & Rodeo sin colision.\\
Objetivo movil & Seguimiento estable.\\
Celda & No invadir zonas de otros equipos.\\
Flota A* & 3 robots completan mision y recarga.\\
\bottomrule
\end{tabularx}

\vspace{0.25cm}
Validacion headless final: 4,328 m $\rightarrow$ 0,739 m, estado ARRIVED, 22 sensores activos, 8 escenarios observados y distancia minima a obstaculo 0,191 m. Flota: 16 planes A*, 8 conflictos evitados, 4 cargas y 3/3 robots READY.
}

\slide{Conclusiones}{
\begin{itemize}
  \item Campos potenciales: solucion sencilla y didactica.
  \item La evitacion reactiva mejora seguridad de navegacion.
  \item La flota A* demuestra coordinacion multi-robot y gestion de bateria.
  \item CoppeliaSim permite ajustar parametros antes de la demostracion.
  \item Limitacion: posibles minimos locales y oscilaciones.
\end{itemize}
}
\end{document}
